\subsubsection{An unconditional first compact comparison for the actual theta series}
\label{sec:firstcertifiedrectangle}
The strict finite-mode hypothesis in
Theorem~\ref{thm:compacttransfer} can be verified
on an explicit, nontrivial rectangle.
Unlike tests of the abstract transfer algebra, this
calculation uses the first two summands of the
actual Jacobi-theta kernel. It certifies a complete
two-variable region without discretizing time
or the Fourier argument.

\begin{theorem}[A certified full-time, first-order rectangle]
\label{thm:firstcompactrectangle}
Let $T=1/2$, $H_{t,2}(x)=
\int_0^\infty e^{tu^2}f_2(u)\cos(xu)\,du$,
where $f_2=\phi_1+\phi_2$ consists of
the first two genuine theta summands.
Write
\[
U_{j,2}(s,x)=L_j[H_{T-s,2}](x),\qquad
K=[0,T]\times[-4,4].
\]
Then, everywhere on $K$,
\begin{equation}\label{eq:firstrectanglefinite}
\boxed{6U_{2,2}(s,x)-U_{1,2}(s,x)<-2\cdot10^{-6}.}
\end{equation}
Let $e_{j,2}$ be the global theta truncation
errors defined in \eqref{eq:compacttransfererror}.
They satisfy
\[
6e_{2,2}+e_{1,2}<10^{-7},
\]
and consequently the required \emph{strict
finite-mode comparison certificate} holds:
\begin{equation}\label{eq:firstrectanglecertified}
\boxed{6U_{2,2}-U_{1,2}
\le -6e_{2,2}-e_{1,2}\quad\text{throughout }K.}
\end{equation}
It follows unconditionally that the
\emph{complete} theta heat flow obeys
\[
\boxed{6U_2(s,x)\le U_1(s,x)
\qquad ((s,x)\in K).}
\]
\end{theorem}

\begin{proof}
For $j\ge0$ let
\[
m_j(t)=\int_0^\infty
u^j e^{tu^2}f_2(u)\,du .
\]
Every integrand is nonnegative, hence
$m_j(t)$ is nondecreasing on $0\le t\le T$.
The exact validated quadrature below proves
the following coarse \emph{strict} bounds,
chosen rationally for easy independent replay:
\begin{align}
m_0(0)&>\frac1{20},
&m_2(0)&>\frac1{2000},\nonumber\\
m_0(T)&<\frac1{10},
&m_1(T)&<\frac1{100},\nonumber\\
m_2(T)&<\frac1{1000},
&m_3(T)&<\frac1{5000},
&m_4(T)&<\frac1{30000}.
\label{eq:rectangle-moment-bounds}
\end{align}

The inequality $\cos y\ge1-y^2/2$,
applied to the positive theta integrand,
shows that for $|x|\le r_0=4$,
\begin{align*}
H_{t,2}(x)&\ge
m_0(0)-\frac{r_0^2}{2}m_2(T)>a,\\
-H_{t,2}''(x)&\ge
m_2(0)-\frac{r_0^2}{2}m_4(T)>b,
\end{align*}
where
\[
a=\frac1{20}-\frac{r_0^2}{2\cdot1000},
\qquad
b=\frac1{2000}-\frac{r_0^2}{2\cdot30000}.
\]
Both $a,b$ are positive. Hence the
first Laguerre expression has the
\emph{uniform} lower bound
\begin{equation}
U_{1,2}=(H_{t,2}')^2-H_{t,2}H_{t,2}''
>ab.
\label{eq:rectangle-U1}
\end{equation}
For the second expression the exact
derivative formula is
\[
U_{2,2}=
\frac{H_{t,2}H_{t,2}^{(4)}}{12}
-\frac{H_{t,2}'H_{t,2}^{(3)}}3
+\frac{(H_{t,2}'')^2}{4}.
\]
Using $|H_{t,2}^{(j)}(x)|\le m_j(T)$,
we obtain
\begin{equation}
U_{2,2}<
b_2:=\frac{1}{10\cdot30000\cdot12}
+\frac{1}{100\cdot5000\cdot3}
+\frac{1}{4\cdot1000^2}.
\label{eq:rectangle-U2}
\end{equation}
The exact rational calculation gives
\[
6b_2-ab=-\frac{79}{30\,000\,000}
<-2\cdot10^{-6},
\]
proving \eqref{eq:firstrectanglefinite}.

For completeness we specify the analytic
trust boundary of \eqref{eq:rectangle-moment-bounds}.
The interval certificate computes the integrals
for $t=0$ and $t=1/2$ over $0\le u\le2$
by outward-rounded composite Simpson quadrature
with respectively $8192$ and $16384$ panels.
The first five derivatives of the theta summands
are controlled by
\[
\phi_m^{(j)}(u)=q_me^{5u-q_me^{4u}}
P_j(q_me^{4u}),\quad
P_{j+1}(v)=(5-4v)P_j(v)+4vP_j'(v).
\]
Writing $B_j=\|P_j\|_1(j+3)!$
gives the valid global bounds
\[
(B_0,B_1,B_2,B_3,B_4)
=(30,1272,79320,6800400,762395760).
\]
For $0\le u\le2$, the derivatives of
$u^j e^{tu^2}$ are bounded by $8$
times the positive-coefficient polynomial
envelope obtained recursively from
\[
R_{j,0}(u)=u^j,\qquad
R_{j,h+1}=R_{j,h}'+uR_{j,h},
\]
evaluated at $u=2$.
Leibniz's rule therefore provides the
explicit fourth-derivative bound
\[
M_{4,j}=
16\sum_{a=0}^4{4\choose a}
B_aR_{j,4-a}(2).
\]
The corresponding Simpson error is at most
\[
\frac2{180}\left(\frac2N\right)^4M_{4,j}.
\]
The omitted integrals $u>2$ are smaller
than $10^{-100}$ for $j\le4$, using
$u^j\le j!e^u$, $e^{tu^2}\le e^{e^{4u}}$,
$3<q_m<16$ and
$e^8>1000$, $e^3>20$.
After $v=e^{4u}$, the tail is at most
\[
j!\,256\int_{1000}^\infty
v^2e^{-2v}\,dv
<10^{-100}.
\]
Every transcendental Simpson node is
enclosed by \texttt{mpmath.iv} 1.3.0,
converted to exact rational endpoints;
the displayed moment bounds are
verified in both separate resolutions.

Finally Theorem~\ref{thm:compacttransfer}
uses the uniform theta heat-tail
coefficient
\[
e_{j,2}
=\frac{3(2j+1)C_*}{2\cdot4^{2j}}\,
q_3e^{-q_3+1/16},\quad q_3=9\pi,\quad C_*<100.
\]
The rational bounds $27<q_3<198/7$,
$e^3>20$, and $q_3-1/16>24$ yield
\[
e_{j,2}<
\frac{3(2j+1)100}{2\cdot4^{2j}}
\frac{198/7}{20^8}.
\]
In particular,
\[
6e_{2,2}+e_{1,2}
<5.05\cdot10^{-8}<10^{-7}.
\]
The exact rational inequalities give even the
stronger bound
\[
6b_2-ab+6e_{2,2}+e_{1,2}<-2\cdot10^{-6},
\]
which proves \eqref{eq:firstrectanglecertified}
without division by a Laguerre expression.
The compact transfer theorem then yields the
claimed genuine-theta inequality.
\end{proof}

This first successful rectangle already spans
the \emph{entire backward heat-time interval}
$s\in[0,1/2]$, including $t=0$.
Its Fourier range is, however, confined to
$|x|\le4$, and its Laguerre order is only
$n=1$. The proof relies mainly on
positivity of the underlying theta moments,
not on a mechanism excluding off-critical
zeros. Extending this result to every order
and unbounded Fourier height remains the
unproved global problem. In particular it
would be incorrect to infer RH from
\eqref{eq:firstrectanglecertified} alone.

The complete interval source is
\path{certificates/xi_first_compact_theta_rectangle.py};
the kernel-checked rational algebra is in
\path{XiKernel/FirstCompactThetaRectangle.lean}.
The integral enclosures and complete analytic
heat/Fourier passage have not been certified
end to end in Lean.


\begin{proposition}[Every fixed order admits a compact rectangle]
\label{prop:allorderslocalrect}
For every integer $n\ge1$ there exist
$r_n>0$, an integer $M_n\ge1$ and
a finite constant $c_n>0$ such that,
with $K_n=[0,1/2]\times[-r_n,r_n]$
and $C_{K_n}=c_n$, the actual
finite theta approximation obeys
\[
k_nU_{n+1,M_n}-c_nU_{n,M_n}
\le-k_ne_{n+1,M_n}-c_ne_{n,M_n}
\qquad\hbox{throughout }K_n.
\]
This statement does not assert a
common positive radius for all orders
or any control at unbounded Fourier height.
\end{proposition}
\begin{proof}
For $H_t(iy)=\int_0^\infty
e^{tu^2}\Phi(u)\cosh(yu)\,du$,
all even Taylor coefficients are positive.
Consequently
$U_n(s,0)=L_n[H_{1/2-s}](0)>0$
for every $n$ and every $s\in[0,1/2]$.
The positive-moment coefficients
at $t=0$ provide a uniform strictly
positive lower bound at $x=0$
throughout the heat slab.
Joint continuity then gives $r_n>0$
and $\eta_n>0$ such that
$U_n(s,x)\ge\eta_n$ on $K_n$.
The next-order expression is bounded
on $K_n$, say $|k_nU_{n+1}|\le A_n$.
Choose
$c_n=(A_n+2)/\eta_n$; then
\[
k_nU_{n+1}-c_nU_n\le-2
\quad\hbox{on }K_n.
\]
The explicit theta heat truncation
errors $e_{j,M}$ tend to zero
as $M\to\infty$ for every fixed order.
Choose $M_n$ so large that
\[
2\bigl(k_ne_{n+1,M_n}
+c_ne_{n,M_n}\bigr)<1.
\]
Since the finite-mode defect differs
from the genuine defect by at most
$k_ne_{n+1,M_n}+c_ne_{n,M_n}$,
the finite-mode defect is at most
$-2+\mathrm{budget}\le-\mathrm{budget}$.
This proves the claim.
\end{proof}

The quantifier order in this proposition
is essential: the radius, mode cutoff
and comparison coefficient may all depend
on $n$. The argument uses positivity at
the central Fourier point, not a mechanism
excluding nonreal zeros. It supplies
\emph{no} bound when $|x|\to\infty$,
which is precisely the unresolved
global step in the Riemann hypothesis programme.
